\subsection{GAN Evaluation}
\label{subsec:gan_evaluation}

To quantitatively evaluate the quality and realism of the synthetic
chest X-ray images generated by the DCGAN, the Fréchet Inception
Distance (FID) \cite{Heusel2017FID} was used as the primary
metric. The FID computes the distance between the feature
distributions of real and synthetic images, where a lower score
designates a higher similarity to the real dataset.

Upon completing the 7000 training epochs, the generative model
achieved a final FID score of $25.37$. This outcome indicates a
successful generation of synthetic data that accurately reflects the
structural properties of genuine patient scans. By demonstrating this
level of visual fidelity, the augmented images were validated for use
in equilibrating the dataset. This successfully augmented the
\texttt{NORMAL} minority class, overcoming the initial data imbalance
and providing a robust foundation for the downstream classification
tasks. Furthermore, this data augmentation framework is highly
extensible and can be adapted to broader medical domains requiring 
dataset equilibration.